\chapter{Market Overview and Regional Viability}

\section{Global and Regional Market Size}

Estimates of the global AMR market for 2025 cluster between roughly USD
2.3~billion and USD 4.9~billion, growing at about 14--16\% per year, with the
spread driven by whether ``AMR'' is counted narrowly (mobile bases) or broadly
(bases plus software and services) \cite{valuates_agv_amr}. Demand is led by
warehousing and manufacturing, with healthcare the fastest-growing end-use.

Adjacent segments give a sense of scale and direction without being directly
addressable by a small Tunisian entrant: the retail-robotics market has been
reported at around USD 20~billion in 2024 with projections toward roughly USD
181~billion by 2031 \cite{gii_retail2030}, and the Middle East \& Africa
industrial-robotics market has been put at about USD 1.45~billion in 2024 rising
toward USD 4.29~billion by 2030 (around 19\% per year) \cite{grandview_mea}.
These figures matter mainly as context: the regional money is real and growing,
but it is concentrated in the GCC and in segments whose economics assume
high labor costs.

\section{Is the Market Viable in Tunisia / North Africa / MENA?}

\subsection{Favorable Factors}

\begin{itemize}
    \item Tunisia hosts one of the largest startup ecosystems in MENA, with well
          over 1,000 active startups and a top-five regional ranking
          \cite{lucidity_tunisia2025,startupgenome_tunisia,startuplist_tunisia}.
    \item The Tunisian Startup Act offers labelled startups tax advantages,
          foreign-currency access and simplified procedures (administered via
          Smart Capital), partially offsetting an otherwise restrictive
          financial regime.
    \item Public policy explicitly encourages robotics and industrial innovation
          \cite{sixwresearch_tunisia,meanceo_funding}.
    \item Local manufacturing of robots is already proven: TIRA Robots is a
          Tunisian maker of industrial robotic arms, and Enova Robotics has
          built and exported a security UGV -- evidence that designing and
          building robots locally is feasible \cite{wamda_tira}.
    \item Deep, trilingual engineering talent (on the order of 10,000 new
          engineers per year) at a competitive wage base.
\end{itemize}

\subsection{Unfavorable / Region-Specific Factors}

\begin{itemize}
    \item A documented brain drain: a large share of Tunisian tech engineers
          work for European companies, partly driven by currency restrictions.
    \item No MENA country has yet built a mature robotics industry; doing so
          requires significant capital, and global leaders already dominate.
    \item Core supply chains (motors, LiDAR, batteries, compute boards) are not
          manufactured locally, so import dependency affects both cost and
          lead time.
    \item There is no Tunisia-specific automation subsidy comparable to the
          GCC's national programmes; EU-funded industrial-modernisation pilots
          are a more realistic co-financing channel.
\end{itemize}

\subsection{Strategic Implication}

The right strategy is \emph{not} to copy an international AMR and sell it
cheaper -- that is structurally impossible while the high-value components stay
imported. The real opportunity is in modular architecture, software, local
support, and a pricing model designed for the local market (rental, sharing
between smaller companies, usage-based pricing). The practical build structure
that follows from this is hybrid: import only the irreducible high-tech
components (sensors, drives, GPUs) under the most favourable customs regimes
available to exporters, and fabricate the chassis, structural modules and
assembly locally -- which also underpins the local support advantage that global
competitors lack.
